% Generated draft: requires mathematical and source review.
\documentclass[11pt]{article}
\usepackage[margin=25mm]{geometry}
\usepackage{fontspec}
\setmainfont{Latin Modern Roman}
\usepackage{amsmath,amssymb,mathtools,unicode-math}
\setmathfont{Latin Modern Math}
\usepackage{enumitem,needspace,xurl,hyperref}
\setlength{\parindent}{0pt}
\setlength{\parskip}{5pt}
\setlength{\emergencystretch}{3em}
\newcommand{\N}{\mathbb{N}}
\newcommand{\Z}{\mathbb{Z}}
\newcommand{\Q}{\mathbb{Q}}
\newcommand{\R}{\mathbb{R}}
\newcommand{\C}{\mathbb{C}}
\newcommand{\F}{\mathbb{F}}
\newcommand{\A}{\mathbb{A}}
\newcommand{\PP}{\mathbb{P}}
\newcommand{\E}{\mathbb{E}}
\newcommand{\eps}{\varepsilon}
\newcommand{\dd}{\,\mathrm{d}}
\begin{document}
\section*{UnsolvedMath ID: 30006581}
\textbf{Source problem code:} catalogue entry
\subsection{Definitions and mathematical statement}
Write \(T_n=(\mathbb{G}_m)^n\) for the \(n\)-dimensional algebraic torus. Using algebraic points as an explicit working interpretation of the supplied background, put
\[
T_n(\overline{\Q})=(\overline{\Q}^{\times})^n,
\qquad
\Sigma_d(T_n)=\bigcup_{\substack{H\leq T_n\text{ algebraic subgroup}\\ \operatorname{codim}_{T_n}H\geq d}}H(\overline{\Q}),
\quad d\in\Z_{\geq1}.
\]
Here the union is over algebraic subgroups, without an added connectedness restriction; it is not a union over all their translates. A translate of \(H\) means a coset \(gH\). Let \(h\) denote a Weil height on torus points. The source does not specify its normalization or projective embedding, so this notation does not determine a numerical explicit bound.

\textbf{Supplied background.} The record attributes the following statement to Bombieri, Masser and Zannier, with source-local reference label [4]. Let \(n\geq2\), and let \(C\subset T_n\) be a geometrically irreducible algebraic curve, defined over \(\Q\) as rendered in the supplied excerpt, which is not contained in a translate of any proper algebraic subgroup of \(T_n\). Then
\[
\exists B_C<\infty\ \ \forall s\in C(\overline{\Q})\cap\Sigma_1(T_n),\qquad h(s)\leq B_C,
\]
and
\[
\#\bigl(C(\overline{\Q})\cap\Sigma_2(T_n)\bigr)<\infty.
\]
The excerpt describes the second assertion as a special case of the Zilber--Pink conjecture. It says that the proof of the first assertion can in principle be made effective and attributes an explicit height bound to the author rendered as ``Habbeger,'' with source-local reference label [8]. These are supplied contextual assertions, not independently verified statements here. The excerpt begins a second theorem but omits its actual height inequality.

\textbf{Requested research target.} Determine fully explicit height bounds obtained using \(G\)-functions for unlikely intersections in algebraic tori and abelian varieties. The supplied clean statement specifically requests extending the stated line-section results in \(\mathbb{G}_m^3\) to algebraic curves in \(\mathbb{G}_m^n\), and developing analogous bounds for curves in \(A^n\), where \(A\) is a simple abelian variety. An abelian variety over a specified field is simple over that field if it has no nonzero proper abelian subvariety defined over that field; geometric simplicity requires this condition after passage to an algebraic closure. The record does not choose between these meanings.

The target is a research program with unresolved formulation parameters. The phrase ``arbitrary algebraic curves'' does not state whether the geometric irreducibility and exclusion of subgroup translates from the background are retained. No precise exceptional locus, subgroup codimension threshold, or meaning of ``very unlikely'' is supplied for the requested extensions. For the abelian setting, the field of definition, polarization, height, simplicity convention, and necessary restrictions on the curve are unspecified. The record also gives neither the admissible input data for a fully explicit bound nor a comparison criterion for ``strongest.'' Consequently, the supplied evidence does not determine a unique quantified bound for either extension. These qualifications must be resolved before replacing this programmatic target by a precise universal inequality.
\subsection{Short English statement}
The problem asks for explicit height bounds using \(G\)-functions, extending supplied results for line sections in \(\mathbb{G}_m^3\) to curves in higher-dimensional tori and to curves in powers of a simple abelian variety. The excerpt supplies toric background but omits the explicit inequality and several hypotheses needed for the extensions. Its \texttt{partially\_solved} status is supplied dated curation, not an independently verified conclusion.
\subsection{Sources}
\begin{itemize}
\item Local UnsolvedMath registry: supplied record with permanent ID \(30006581\), source problem code \texttt{catalogue entry}, title ``Effective G-Function Bounds for Unlikely Intersections.'' Project registry paths are \texttt{data/unsolvedmath/problems.json} and the supplemental \texttt{catalog/problems.json}; this draft uses only the supplied record.
\item Supplied report citation: \emph{Model Theory and Arithmetic} (2026), Oberwolfach Reports, DOI \texttt{10.4171/owr/2026/3}, pp.~157--216. Supplied literature-source label: ``Oberwolfach Report 3/2026.'' \url{https://doi.org/10.4171/owr/2026/3}. The excerpt attributes the contribution ``Explicit height bounds for (very) unlikely intersections using G-functions'' to Martin Orr; the separate proposer field is empty.
\item Supplied literature reference: Orr, \emph{Explicit height bounds for G-functions and unlikely intersections with lines in tori}. \url{https://arxiv.org/abs/2607.09565}. The record describes this as a July 2026 preprint.
\item Contextual attributions retained from the excerpt: Bombieri, Masser and Zannier [4] for its Theorem 1, and ``Habbeger'' [8] for the explicit bound mentioned before the truncated Theorem 2. The bibliographic entries corresponding to [4] and [8] are not supplied.
\item Literature assessment date: \texttt{2026-08-22}. This is supplied dated curation, not independently verified. It labels the problem \texttt{partially\_solved}, states that explicit height bounds are known in the classical torus setting, and credits Orr's July 2026 preprint with explicit \(G\)-function bounds and applications to lines in tori while describing the broader very-unlikely-intersection program as unresolved.
\item The record labels its statement \texttt{corrected\_verified} and reports that the official report was checked directly and that OCR and notation were repaired. That verification claim belongs to the supplied record; no independent source visit or verification is asserted here.
\end{itemize}
\end{document}
